\uuid{RYXc}
\titre{Loi de la moyenne d'un échantillon normal}
\auteur{Maxime Nguyen}
\organisation{AMSCC}
\datecreate{2026-11-27}
\niveau{L2}
\module{Probabilité}
\chapitre{Probabilité continue}
\sousChapitre{Loi normale}
\difficulte{1}
\theme{loi normale, moyenne empirique, somme de variables indépendantes}

\contenu{
\begin{question}
\qtype{onlyone}
On note $\mathcal{N}(\mu,\sigma)$ la loi normale d'espérance $\mu$ et d'écart-type $\sigma$. La taille (en cm) d'un adulte est modélisée par la loi $\mathcal{N}(170, 10)$. On mesure $25$ adultes ; leurs tailles $X_1, \dots, X_{25}$ sont indépendantes, de même loi $\mathcal{N}(170, 10)$. La moyenne $\overline{X} = \frac{1}{25}\sum_{i=1}^{25} X_i$ suit la loi :
\begin{answers}
\bad{$\mathcal{N}(170, 10)$}
\bad{$\mathcal{N}(170, 4)$}
\bad{$\mathcal{N}(170;\, 0{,}4)$}
\good{$\mathcal{N}(170, 2)$}
\end{answers}
\begin{explanations}
La somme suit $\mathcal{N}(25 \times 170,\ \sqrt{25} \times 10)$,
soit $\mathcal{N}(4250, 50)$ ; en divisant par $25$, on obtient
$\mathcal{N}\left(170, \frac{50}{25}\right) = \mathcal{N}(170, 2)$. En général,
$\overline{X} \sim \mathcal{N}\left(\mu, \frac{\sigma}{\sqrt n}\right)$.

A : la moyenne de $25$ mesures est moins dispersée qu'une mesure seule.

B : $\frac{\sigma^2}{n} = 4$ est la variance de $\overline{X}$.

C : $\frac{\sigma}{n}$, on a divisé par $n$ au lieu de $\sqrt n$.

Contre-sens visés : croire que moyenner ne réduit pas la dispersion (A) ;
diviser l'écart-type par $n$ au lieu de $\sqrt n$ (C) ; confondre écart-type et
variance (B).
\end{explanations}
\end{question}
}
